\documentclass[12pt, letterpaper]{article}
\usepackage[margin=1.2in]{geometry}
\usepackage{ebgaramond}
\usepackage[T1]{fontenc}
\usepackage{microtype}
\usepackage{setspace}
\usepackage{parskip}
\usepackage{titling}
\usepackage{fancyhdr}
\usepackage[utf8]{inputenc}

\setlength{\parindent}{2em}
\onehalfspacing

\pagestyle{fancy}
\fancyhf{}
\fancyhead[L]{\small\textit{The Bell}}
\fancyhead[R]{\small\textit{NASA Mission Series v1.20}}
\fancyfoot[C]{\small\thepage}
\renewcommand{\headrulewidth}{0.4pt}

\title{\Huge\textbf{The Bell}\\[0.5em]\large\textit{A story from Mission 1.20 --- Mercury-Redstone 4: Liberty Bell 7}}
\author{NASA Mission Series\\[0.3em]\small Miles Tiberius Foxglove --- Mukilteo, WA}
\date{March 30, 2026}

\begin{document}
\maketitle
\thispagestyle{empty}

\vspace{1em}
\noindent\rule{\textwidth}{0.4pt}
\vspace{1em}

\noindent The capsule was floating. That was the first clear fact. Gus Grissom lay on his back in the couch, the harness still tight across his chest, and the capsule was floating on the Atlantic Ocean and everything had gone right. Fifteen minutes and thirty-seven seconds from liftoff to splashdown, a suborbital arc nearly identical to Shepard's, the instruments reading nominal across the board, the reentry load peaking at eleven point one g --- slightly less than Shepard had endured, and Grissom had stayed conscious through all of it, watching the gauges, hand on the controller, the backup system that was always the pilot.

The capsule bobbed. He could hear waves slapping the hull, the irregular rhythm of the Atlantic. The porthole showed blue sky, a wisp of cloud, the bright flatness of mid-morning ocean. He was home. Second American in space. July 21, 1961.

He began the post-landing checklist. Remove helmet neck seal. Disconnect oxygen hose. Unfasten harness --- leave the lap belt connected until the helicopter arrived. Record cockpit switch positions. He pulled the safety pin on the hatch detonator and placed it in the pocket of his suit where it would be found later. The hatch was armed. He would blow it on command, after the helicopter attached the recovery cable, after the capsule was safely hooked to the horse collar. Standard procedure. Everything in order.

He called the recovery helicopter on the radio. ``Give me how much longer it's going to be before you get here.''

The helicopter was four minutes out. He waited. The capsule rocked gently. Through the porthole, the sky was very blue.

\begin{center}
$\bullet\quad\bullet\quad\bullet$
\end{center}

The sound was like a shotgun fired inside a tin can. Seventy explosive bolts detonating simultaneously --- the Mild Detonating Fuse ring that held the hatch cover to the capsule frame severing in one instant, a single percussive crack that compressed the air inside the capsule and then released it in a rush as the hatch cover flew outward and the Atlantic Ocean poured in.

Grissom did not fire the hatch. He would say this for the rest of his life, to anyone who asked, in debriefings and press conferences and private conversations that grew more bitter as the years passed. He did not hit the plunger. The plunger required a deliberate strike with the heel of the palm --- hard enough to leave a bruise. When they examined his hands afterward, there was no bruise. The plunger cap was in place, undisplaced. He had been lying there, minding his own business, waiting for the helicopter, and the hatch had blown on its own.

But that was later, the analysis and the accusation. In the moment, there was only water.

Seawater came through the open hatch like a river. The opening was nineteen inches in diameter --- large enough to crawl through, which was the point, which was what he needed to do right now. The capsule was taking on water at a rate that would fill it in under two minutes. Grissom disconnected the lap belt, grabbed the instrument panel, and pulled himself through the hatch opening into the ocean.

\begin{center}
$\bullet\quad\bullet\quad\bullet$
\end{center}

\emph{The Giant Pacific Octopus makes its den in rock crevices on the floor of Puget Sound. Enteroctopus dofleini: the largest species of octopus on Earth, arms spanning six meters, body mass up to fifty kilograms. It is intelligent enough to solve mazes, open jars, recognize individual human faces. It lives two to five years. In that time, a female will lay between twenty thousand and one hundred thousand eggs in a single clutch, suspended in braided strands from the ceiling of her den, and she will tend them for six to seven months without eating. She ventilates the eggs with jets from her siphon. She cleans them with the tips of her arms. She does not leave the den. When the eggs hatch, she dies.}

\emph{The octopus is an animal of extraordinary investment. Everything it has goes into one act. The den is chosen with care --- a crevice the right size, in the right depth, with the right current flow. The eggs are laid with precision. The brooding is absolute. And when it is done, the animal is spent. There is no second clutch. There is no recovery. The octopus gives everything it has to one moment, and then the moment passes, and the tide carries what remains.}

\begin{center}
$\bullet\quad\bullet\quad\bullet$
\end{center}

Grissom was in the water. The suit was not designed for swimming. It was a pressure suit --- a modified Navy Mark IV, bulky and stiff, with a neck seal that was supposed to keep water out. But the suit had an open oxygen inlet port at the chest. He hadn't remembered to close it. Water began seeping in through the port, slowly at first, then faster as each wave washed over him. The suit was filling. The boots were heavy. The gloves made it difficult to tread water.

He could see the capsule. Liberty Bell 7 was listing badly now, the open hatch below the waterline, green dye marker spreading in the water around it like blood from a wound. The recovery helicopter was overhead --- he could hear the rotors, the deep thump of the main blades, the whine of the turbine. But the helicopter was going for the capsule, not for him. The pilot, Marine Lieutenant Jim Lewis, had hooked the recovery cable to the capsule's neck and was trying to lift.

Grissom treaded water and watched his spacecraft die.

The capsule weighed twenty-four hundred pounds dry. With seawater flooding through the open hatch, the weight was climbing by the second --- a hundred kilograms every few seconds, the Bernoulli equation running in real time through a nineteen-inch orifice. The helicopter strained. The cable went taut. The engine temperature climbed. Jim Lewis held the hover, watching the gauges, the collective stick at maximum pitch, the airframe shuddering as the UH-34D's engine tried to lift a weight that was growing faster than physics allowed.

The engine chip warning light came on. Metal particles in the oil --- the engine was tearing itself apart. Lewis held on. Five seconds. Ten. The capsule didn't come up. It hung in the water, half submerged, the hatch pouring ocean, the weight climbing past the helicopter's maximum hover capacity. Lewis held on another three seconds and then released the cable.

Liberty Bell 7 sank.

\begin{center}
$\bullet\quad\bullet\quad\bullet$
\end{center}

In the water, Grissom was drowning. Not dramatically --- not flailing, not panicking. He was a test pilot. He managed the situation. He kept his head above the surface. He signaled the second helicopter. But the suit was heavy with water now, pulling him down, and each wave pushed his face under, and the helicopter seemed to be circling rather than approaching, and he could feel the weight of the ocean pulling at his boots, and for the first time in the flight he was genuinely afraid.

The second helicopter lowered a horse collar. Grissom grabbed it, pulled it under his arms, and was lifted out of the Atlantic. He hung from the helicopter, water pouring from his suit, spinning slowly as the rotor wash buffeted him, and he looked down and saw the circle of green dye where the capsule had been, and the circle of darker water where the capsule was going, and then the helicopter gained altitude and the circle shrank to a point and then to nothing.

He was alive. The capsule was gone.

\begin{center}
$\bullet\quad\bullet\quad\bullet$
\end{center}

\emph{The Giant Pacific Octopus hunts by ambush. It matches the color and texture of any surface it touches --- rock, sand, kelp, coral --- in under a second, using chromatophores in its skin: tens of thousands of tiny sacs of pigment, each controlled by a ring of muscle, each capable of expanding to display color or contracting to hide it. The octopus does not decide to change color. The chromatophores respond to the texture under the arms, a feedback loop that operates below conscious control. The animal becomes what it touches.}

\emph{But camouflage is not the octopus's only tool. When threatened, it can jet: a contraction of the mantle cavity that forces water through the siphon at high speed, propelling the animal backward in a burst that can reach twenty-five miles per hour. And if jetting is not enough, there is ink --- a cloud of melanin ejected into the water, a decoy that confuses the predator while the octopus escapes. The ink cloud hangs in the water, roughly the size and shape of the octopus itself, a ghost of the animal that was there a moment ago.}

\emph{Three defenses: hiding, fleeing, and leaving something behind that looks like you. The octopus uses all three, in order, each one a retreat from the one before. First it hides. When hiding fails, it runs. When running is not enough, it leaves a ghost of itself to take the blow.}

\begin{center}
$\bullet\quad\bullet\quad\bullet$
\end{center}

Grissom spent the rest of his life defending himself. NASA held a review and cleared him --- the evidence supported his account, the physical findings were consistent with an uncommanded hatch blow --- but the public narrative had already formed. The astronaut panicked. He blew the hatch. He lost the spacecraft. The accusation followed him through Gemini, through the early Apollo program, through the years of training and preparation that led to January 27, 1967, when Grissom and Roger Chaffee and Ed White were strapped into the Apollo 1 capsule on the launch pad at Cape Kennedy for a plugs-out test, and a fire started in the pure oxygen atmosphere, and the hatch --- a hatch that opened inward, a hatch that could not be opened against internal pressure, a hatch that had been redesigned specifically because of what happened to Liberty Bell 7 --- could not be opened in time.

Grissom died at thirty-seven. He died in a spacecraft, which was the one place he was not afraid to be. The hatch that killed him in Apollo 1 was the descendant of the hatch that had tried to kill him in Mercury. The redesign had gone the wrong direction: afraid of another accidental opening, they had built a hatch that could not be opened at all. The lesson of Liberty Bell 7 --- that hatches should open easily in emergencies --- had been inverted by the fear that Grissom had opened one when he shouldn't have. The accusation that haunted him became, in a terrible and literal way, the thing that killed him.

After the fire, they redesigned the hatch again. Apollo spacecraft from Block II onward had a quick-release hatch that could be opened in seven seconds. It was the hatch that carried men to the Moon. It was the hatch that should have been there on January 27, 1967. It was the hatch that Grissom would have wanted.

\begin{center}
$\bullet\quad\bullet\quad\bullet$
\end{center}

Ernest Hemingway wrote \emph{The Old Man and the Sea} in 1951, ten years before Liberty Bell 7. Santiago, the old fisherman, hooks the great marlin far out in the Gulf Stream and fights it for three days. He lashes it to the side of his skiff and turns for home. The sharks come. Santiago fights them with his harpoon, then his knife, then a club made from the broken tiller. One by one, the sharks strip the marlin to the bone. By the time Santiago reaches the harbor, there is nothing left of the fish but the skeleton lashed to the hull. He has won nothing. He has lost everything. He climbs from the skiff, carries the mast to his shack, and falls asleep on his cot.

``A man can be destroyed but not defeated.''

Hemingway died on July 2, 1961 --- nineteen days before Grissom's flight. He had been ill for years. The electroshock treatments at the Mayo Clinic had damaged his memory, which for a writer was the cruelest possible injury. He could no longer write. He could no longer be the thing he was. He chose his own ending, in Ketchum, Idaho, with a shotgun he had owned for years, and the morning was quiet and the mountains were there and the trout were still running in Silver Creek, and the work he had done remained on the shelves even though the man who had done it did not remain.

Santiago was beaten but not defeated. Hemingway was both. Grissom was neither --- he was accused of something he did not do, and the accusation cost him more than the ocean had taken, and he kept flying anyway, because that was what a pilot did.

\begin{center}
$\bullet\quad\bullet\quad\bullet$
\end{center}

\emph{When the female Giant Pacific Octopus dies after her eggs hatch, her body drifts to the seafloor. Within hours, scavengers arrive --- crabs, sea stars, smaller octopuses. The tissue is consumed. The beak, made of chitin, persists longer, but eventually it too is broken down. Nothing remains of the animal except the twenty thousand to one hundred thousand hatchlings, each one smaller than a grain of rice, drifting in the plankton layer, transparent and fragile and carrying the full genetic code of a creature that could solve mazes and open jars and change color faster than thought.}

\emph{The hatchlings are planktonic for four to twelve weeks. Most will be eaten. The survival rate is estimated at one to two percent. Of a clutch of fifty thousand, perhaps five hundred to a thousand will survive to settle on the bottom and begin their own two-to-five-year lives. This is not tragedy. It is investment. The octopus does not mourn the ninety-eight percent that are lost. It produced a surplus because the ocean is large and the risks are many, and the species continues because the investment, even at two percent return, is sufficient.}

\begin{center}
$\bullet\quad\bullet\quad\bullet$
\end{center}

Liberty Bell 7 sank to four thousand nine hundred meters. Nearly five kilometers of water. The pressure at that depth is four hundred and eighty-seven atmospheres --- every square centimeter of the capsule's hull bearing fifty kilograms of force. The capsule survived because the hatch was open: water had flooded in during the descent, equalizing the pressure inside and out. The hull was never subjected to a crushing differential. It simply sank, tumbling slowly in the darkness, the open hatch admitting the ocean at every depth, until it reached the abyssal plain of the Atlantic and settled on the sediment and the current stopped and the silence was absolute.

For thirty-eight years, the capsule lay in darkness. No sunlight reaches four thousand nine hundred meters. The abyssal plain is a desert of fine sediment, three to five degrees Celsius, perpetually dark, sparsely populated by organisms adapted to an environment that would crush a submarine. The capsule was beyond reach of any technology available in 1961. It was, for all practical purposes, gone.

Curt Newport did not accept this. A former Navy diver and salvage specialist, Newport spent fourteen years searching for Liberty Bell 7. He used historical records to narrow the search area, side-scan sonar to map the seafloor, and eventually a remotely operated vehicle to identify and photograph the capsule. On July 20, 1999 --- thirty-eight years and one day after the flight, and thirty years to the day after Apollo 11 landed on the Moon --- a recovery ship lowered cables to the abyssal plain and attached them to Liberty Bell 7.

The lift took five hours. The capsule rose through the zones of the ocean in reverse: abyssal to midnight to twilight to photic, the water around it brightening imperceptibly, then faintly, then unmistakably as the sunlight that had not touched it in thirty-eight years found it again. When it broke the surface, seawater poured from the open hatch --- the same hatch, the same opening, still admitting the ocean, still the wound that had started everything.

Inside the capsule, they found Grissom's equipment. His survival knife. His mirrors and dye markers. The instrument panel, corroded but readable. And the safety pin from the hatch detonator, still in the pocket of the suit where Grissom had placed it forty minutes into the flight, before the world changed.

\begin{center}
$\bullet\quad\bullet\quad\bullet$
\end{center}

Hemingway again: ``The old man was now definitely and finally \emph{salao}, which is the worst form of unlucky.'' Santiago had gone eighty-four days without a fish. The other fishermen thought he was finished. The boy's parents wouldn't let the boy sail with him anymore. The old man went out alone on the eighty-fifth day and hooked the great fish and brought it home destroyed.

Grissom went out on his eighty-fifth day. He flew Gemini 3 in 1965 --- the first crewed Gemini mission --- and named the capsule \emph{Molly Brown}, after the ``unsinkable'' Molly Brown of Titanic fame. It was a joke. It was also not a joke. NASA told him the name was undignified and asked him to change it. He offered ``Titanic'' as an alternative. NASA allowed \emph{Molly Brown}. After that, they stopped letting astronauts name their spacecraft. The tradition didn't resume until Apollo 9.

Grissom flew \emph{Molly Brown} perfectly. Three orbits, manual reentry, precise splashdown. The capsule floated. The hatch stayed closed. He rode the horse collar to the carrier deck and walked across it without looking back at the water.

\begin{center}
$\bullet\quad\bullet\quad\bullet$
\end{center}

\emph{The Giant Pacific Octopus has three hearts. Two branchial hearts pump blood through the gills. One systemic heart circulates blood through the body. The blood is copper-based --- hemocyanin instead of hemoglobin --- which makes it blue. The blue blood is efficient at carrying oxygen in cold, low-oxygen water. It is less efficient in warm water, which is why the octopus lives in the cold North Pacific, in Puget Sound and the Strait of Juan de Fuca and the kelp forests of the San Juan Islands, where the water temperature rarely exceeds twelve degrees Celsius.}

\emph{Three hearts, blue blood, eight arms, a brain distributed through its entire body. The octopus is an alien architecture --- a body plan so different from ours that convergent intelligence is the only explanation for its problem-solving ability. It did not inherit intelligence from a common ancestor with mammals. It invented it independently, in the ocean, half a billion years ago. The octopus is proof that the universe produces minds from any material available, in any shape that works, and that the shape does not need to be ours.}

\begin{center}
$\bullet\quad\bullet\quad\bullet$
\end{center}

The bell is a symbol. The Liberty Bell --- the real one, in Philadelphia --- cracked. The crack is the most famous thing about it. It cracked in 1835, or 1846, or at some other point depending on which account you trust, and after the crack it could no longer ring, and after it could no longer ring it became more famous than it had ever been when it was whole. The crack made it real. A bell that works perfectly is just a bell. A bell with a crack in it is a story.

Grissom painted the Liberty Bell crack on his capsule. He chose the name. He chose to mark the spacecraft with a symbol of something beautiful and broken, something that cracked and was more famous for cracking than for ringing. Maybe he knew. Not in the specific, prophetic sense --- not that the hatch would blow and the capsule would sink --- but in the way a test pilot knows that the machines he flies are always one failure away from killing him. The crack was already there. It was always there. He painted it on the outside because that's where it belonged --- visible, acknowledged, part of the design.

``The old man was beaten now, he knew that truly.'' Santiago turning for home with the skeleton of the marlin lashed to his hull. ``He settled against the wood comfortably and took his suffering as it came and the fish hung half over the side of the skiff.'' The old man was beaten. The old man was not defeated. The distinction is everything.

Grissom was beaten by the hatch. He was beaten by the ocean. He was beaten by the accusation that followed him for six years until the fire took him out of the conversation entirely. He was not defeated. He flew again. He named his next spacecraft after an unsinkable woman. He was assigned to command the first Apollo mission, which would have put him first in line for a lunar landing. He kept flying because the alternative was not flying, and not flying was not something he could do.

The capsule came back. Thirty-eight years on the bottom of the Atlantic and Curt Newport brought it home. The hatch was still open. The instruments were still readable. The safety pin was still in the pocket. The evidence supported what Grissom had always said: he did not fire the hatch. The bolt circuit had a vulnerability to stray current and salt water intrusion. The bolts had fired themselves.

A man can be destroyed but not defeated. A bell can crack and still be a bell. A capsule can sink to five kilometers and still come home.

\vspace{2em}
\begin{center}
\rule{3cm}{0.4pt}\\[1em]
\textit{Dedicated to Ernest Hemingway (1899--1961)\\who wrote about old men, and the sea, and the difference\\between being beaten and being defeated}
\end{center}

\vspace{2em}
\begin{center}
\small
NASA Mission Series v1.20 --- Mercury-Redstone 4 --- July 21, 1961\\
Organism: \textit{Enteroctopus dofleini} (Giant Pacific Octopus) $\bullet$ Bird: Harlequin Duck (degree 20)
\end{center}

\end{document}
